% Invented example for the complete-paper skill (parts/supplement.md).
% Same preamble, same notation and same citation forms as main.tex.
\documentclass[10pt]{article}
\usepackage[margin=22mm]{geometry}
% One preamble for main.tex and supplement.tex: the same packages, the same notation, the same macros.
% Invented example for the complete-paper skill.
\usepackage{amsmath}
\usepackage{graphicx}
\usepackage{xr}
\newcommand{\jc}{j_{\mathrm{c}}}        % threshold current density
\newcommand{\jzero}{j_{0}}              % threshold at 90 degrees of tilt
\newcommand{\mz}{m_z}                   % out-of-plane component of the octupole moment
\newcommand{\tsh}{\theta_{\mathrm{SH}}} % spin Hall angle
\newcommand{\Ku}{K}                     % anisotropy energy density
\newcommand{\un}[1]{\,\mathrm{#1}}      % units: upright, thin space

\externaldocument[M-]{main}
\renewcommand{\thesection}{S\arabic{section}}
\renewcommand{\thefigure}{S\arabic{figure}}
\renewcommand{\thetable}{S\arabic{table}}
\renewcommand{\theequation}{S\arabic{equation}}

\title{Supplemental Material for:\\ Field-free switching of a tilted kagome antiferromagnet}
\author{A. Author \and B. Author}
\date{}

\begin{document}
\maketitle

\noindent Contents: Sec.~\ref{sec:derivation}, derivation of the threshold, Eq.~(\ref{M-eq:threshold}) of
the main text. Sec.~\ref{sec:robust}, robustness of the threshold. Sec.~\ref{sec:thermal}, effect of
temperature. The data behind every figure, and all runs that are not shown, are in the data package
named in the Data availability statement of the main text.

\section{Derivation of the threshold current}
\label{sec:derivation}

This section derives Eq.~(\ref{M-eq:threshold}) of the main text. The result is Eq.~(\ref{eq:result}).
A reader who accepts the torque balance of Eq.~(\ref{eq:balance}) can skip the rest.

In the frame of the kagome plane the octupole moment is described by one angle $\varphi$, and its energy
density is
\begin{equation}
  E(\varphi) = -\Ku \cos 2\varphi .
  \label{eq:energy}
\end{equation}
The spin current is polarized along $y$. For a plane tilted by $\theta$ about $y$, the damping-like torque
that drives $\varphi$ is the projection
\begin{equation}
  \tau_{\mathrm{DL}} = \frac{\hbar\,\tsh\, j}{2 e d}\,\sin\theta .
  \label{eq:torque}
\end{equation}
The largest restoring torque of Eq.~(\ref{eq:energy}) is $2\Ku$. The moment leaves its energy minimum when
\begin{equation}
  \tau_{\mathrm{DL}} > 2\Ku \quad\text{at}\quad \varphi = \pi/4 ,
  \label{eq:balance}
\end{equation}
and inserting Eq.~(\ref{eq:torque}) gives
\begin{equation}
  \jc(\theta) = \frac{4 e \Ku d}{\hbar\,\tsh}\,\frac{1}{\sin\theta} .
  \label{eq:result}
\end{equation}
The main text absorbs the factor of two into the definition of the anisotropy constant used there.

\section{Robustness of the threshold}
\label{sec:robust}

The main text states that the threshold of Fig.~\ref{M-fig:threshold} moves by less than $5\%$ when
numerical and material parameters are varied. This section gives the evidence. The largest change is
$4.6\%$.

Table~\ref{tab:ranges} lists the reference value of each parameter, the range over which it was varied,
and the largest change of the threshold at $\theta = 30^\circ$. Figure~\ref{fig:robust}(a) shows the
convergence with the cell size. Below $2\un{nm}$ the threshold changes by less than $1\%$, and all runs of
the paper use $1\un{nm}$. Figure~\ref{fig:robust}(b) shows the dependence on damping, the parameter with
the largest effect.

\begin{table}
  \centering
  \caption{Parameters of the simulations, each with its reference value, its range of variation, and the
  largest change of the threshold current at $\theta = 30^\circ$ within that range.}
  \label{tab:ranges}
  \begin{tabular}{lccc}
    \hline
    Parameter & Reference & Range & Change of $\jc$ \\
    \hline
    Damping constant & 0.003 & 0.001 to 0.01 & $4.6\%$ \\
    Cell size & $1\un{nm}$ & 0.5 to $4\un{nm}$ & $0.8\%$ below $2\un{nm}$ \\
    Rise time of the pulse & $10\un{ps}$ & 1 to $50\un{ps}$ & $1.9\%$ \\
    Film thickness $d$ & $20\un{nm}$ & fixed & not varied \\
    \hline
  \end{tabular}
\end{table}

\begin{figure}
  \centering
  \includegraphics[width=0.62\textwidth]{example-image-a}
  \caption{Threshold current density at $\theta = 30^\circ$, relative to the reference run, (a) against
  cell size and (b) against the damping constant. The dashed lines mark a change of $\pm 5\%$.}
  \label{fig:robust}
\end{figure}

The sweep comprises 212 runs. Only the runs that set the entries of Table~\ref{tab:ranges} are drawn in
Fig.~\ref{fig:robust}. The output of every run, with the script that extracts the threshold, is in the
folder \texttt{sweep/} of the data package.

\section{Effect of temperature}
\label{sec:thermal}

The main text states that thermal fluctuations at $300\un{K}$ broaden the transition and do not shift it.
Figure~\ref{fig:thermal} shows the switching probability against current density for 64 runs per point,
each with a different seed of the thermal field. The probability rises from 0.1 to 0.9 over $8\%$ of the
zero-temperature threshold, and it crosses 0.5 within $1\%$ of that threshold.

\begin{figure}
  \centering
  \includegraphics[width=0.45\textwidth]{example-image-b}
  \caption{Switching probability at $300\un{K}$ and $\theta = 30^\circ$ against current density in units
  of the zero-temperature threshold $\jc$. Each point is the mean of 64 runs, and the line is an error
  function fitted to the points.}
  \label{fig:thermal}
\end{figure}

\end{document}
